% Options for packages loaded elsewhere
\PassOptionsToPackage{unicode}{hyperref}
\PassOptionsToPackage{hyphens}{url}
\documentclass[
]{article}
\usepackage{xcolor}
\usepackage[margin=1in]{geometry}
\usepackage{amsmath,amssymb}
\setcounter{secnumdepth}{-\maxdimen} % remove section numbering
\usepackage{iftex}
\ifPDFTeX
  \usepackage[T1]{fontenc}
  \usepackage[utf8]{inputenc}
  \usepackage{textcomp} % provide euro and other symbols
\else % if luatex or xetex
  \usepackage{unicode-math} % this also loads fontspec
  \defaultfontfeatures{Scale=MatchLowercase}
  \defaultfontfeatures[\rmfamily]{Ligatures=TeX,Scale=1}
\fi
\usepackage{lmodern}
\ifPDFTeX\else
  % xetex/luatex font selection
\fi
% Use upquote if available, for straight quotes in verbatim environments
\IfFileExists{upquote.sty}{\usepackage{upquote}}{}
\IfFileExists{microtype.sty}{% use microtype if available
  \usepackage[]{microtype}
  \UseMicrotypeSet[protrusion]{basicmath} % disable protrusion for tt fonts
}{}
\makeatletter
\@ifundefined{KOMAClassName}{% if non-KOMA class
  \IfFileExists{parskip.sty}{%
    \usepackage{parskip}
  }{% else
    \setlength{\parindent}{0pt}
    \setlength{\parskip}{6pt plus 2pt minus 1pt}}
}{% if KOMA class
  \KOMAoptions{parskip=half}}
\makeatother
\usepackage{color}
\usepackage{fancyvrb}
\newcommand{\VerbBar}{|}
\newcommand{\VERB}{\Verb[commandchars=\\\{\}]}
\DefineVerbatimEnvironment{Highlighting}{Verbatim}{commandchars=\\\{\}}
% Add ',fontsize=\small' for more characters per line
\usepackage{framed}
\definecolor{shadecolor}{RGB}{248,248,248}
\newenvironment{Shaded}{\begin{snugshade}}{\end{snugshade}}
\newcommand{\AlertTok}[1]{\textcolor[rgb]{0.94,0.16,0.16}{#1}}
\newcommand{\AnnotationTok}[1]{\textcolor[rgb]{0.56,0.35,0.01}{\textbf{\textit{#1}}}}
\newcommand{\AttributeTok}[1]{\textcolor[rgb]{0.13,0.29,0.53}{#1}}
\newcommand{\BaseNTok}[1]{\textcolor[rgb]{0.00,0.00,0.81}{#1}}
\newcommand{\BuiltInTok}[1]{#1}
\newcommand{\CharTok}[1]{\textcolor[rgb]{0.31,0.60,0.02}{#1}}
\newcommand{\CommentTok}[1]{\textcolor[rgb]{0.56,0.35,0.01}{\textit{#1}}}
\newcommand{\CommentVarTok}[1]{\textcolor[rgb]{0.56,0.35,0.01}{\textbf{\textit{#1}}}}
\newcommand{\ConstantTok}[1]{\textcolor[rgb]{0.56,0.35,0.01}{#1}}
\newcommand{\ControlFlowTok}[1]{\textcolor[rgb]{0.13,0.29,0.53}{\textbf{#1}}}
\newcommand{\DataTypeTok}[1]{\textcolor[rgb]{0.13,0.29,0.53}{#1}}
\newcommand{\DecValTok}[1]{\textcolor[rgb]{0.00,0.00,0.81}{#1}}
\newcommand{\DocumentationTok}[1]{\textcolor[rgb]{0.56,0.35,0.01}{\textbf{\textit{#1}}}}
\newcommand{\ErrorTok}[1]{\textcolor[rgb]{0.64,0.00,0.00}{\textbf{#1}}}
\newcommand{\ExtensionTok}[1]{#1}
\newcommand{\FloatTok}[1]{\textcolor[rgb]{0.00,0.00,0.81}{#1}}
\newcommand{\FunctionTok}[1]{\textcolor[rgb]{0.13,0.29,0.53}{\textbf{#1}}}
\newcommand{\ImportTok}[1]{#1}
\newcommand{\InformationTok}[1]{\textcolor[rgb]{0.56,0.35,0.01}{\textbf{\textit{#1}}}}
\newcommand{\KeywordTok}[1]{\textcolor[rgb]{0.13,0.29,0.53}{\textbf{#1}}}
\newcommand{\NormalTok}[1]{#1}
\newcommand{\OperatorTok}[1]{\textcolor[rgb]{0.81,0.36,0.00}{\textbf{#1}}}
\newcommand{\OtherTok}[1]{\textcolor[rgb]{0.56,0.35,0.01}{#1}}
\newcommand{\PreprocessorTok}[1]{\textcolor[rgb]{0.56,0.35,0.01}{\textit{#1}}}
\newcommand{\RegionMarkerTok}[1]{#1}
\newcommand{\SpecialCharTok}[1]{\textcolor[rgb]{0.81,0.36,0.00}{\textbf{#1}}}
\newcommand{\SpecialStringTok}[1]{\textcolor[rgb]{0.31,0.60,0.02}{#1}}
\newcommand{\StringTok}[1]{\textcolor[rgb]{0.31,0.60,0.02}{#1}}
\newcommand{\VariableTok}[1]{\textcolor[rgb]{0.00,0.00,0.00}{#1}}
\newcommand{\VerbatimStringTok}[1]{\textcolor[rgb]{0.31,0.60,0.02}{#1}}
\newcommand{\WarningTok}[1]{\textcolor[rgb]{0.56,0.35,0.01}{\textbf{\textit{#1}}}}
\usepackage{longtable,booktabs,array}
\newcounter{none} % for unnumbered tables
\usepackage{calc} % for calculating minipage widths
% Correct order of tables after \paragraph or \subparagraph
\usepackage{etoolbox}
\makeatletter
\patchcmd\longtable{\par}{\if@noskipsec\mbox{}\fi\par}{}{}
\makeatother
% Allow footnotes in longtable head/foot
\IfFileExists{footnotehyper.sty}{\usepackage{footnotehyper}}{\usepackage{footnote}}
\makesavenoteenv{longtable}
\usepackage{graphicx}
\makeatletter
\newsavebox\pandoc@box
\newcommand*\pandocbounded[1]{% scales image to fit in text height/width
  \sbox\pandoc@box{#1}%
  \Gscale@div\@tempa{\textheight}{\dimexpr\ht\pandoc@box+\dp\pandoc@box\relax}%
  \Gscale@div\@tempb{\linewidth}{\wd\pandoc@box}%
  \ifdim\@tempb\p@<\@tempa\p@\let\@tempa\@tempb\fi% select the smaller of both
  \ifdim\@tempa\p@<\p@\scalebox{\@tempa}{\usebox\pandoc@box}%
  \else\usebox{\pandoc@box}%
  \fi%
}
% Set default figure placement to htbp
\def\fps@figure{htbp}
\makeatother
\setlength{\emergencystretch}{3em} % prevent overfull lines
\providecommand{\tightlist}{%
  \setlength{\itemsep}{0pt}\setlength{\parskip}{0pt}}
\usepackage{bookmark}
\IfFileExists{xurl.sty}{\usepackage{xurl}}{} % add URL line breaks if available
\urlstyle{same}
\hypersetup{
  pdftitle={Example Report},
  pdfauthor={Magdalena A. Tkacz},
  hidelinks,
  pdfcreator={LaTeX via pandoc}}

\title{Example Report}
\usepackage{etoolbox}
\makeatletter
\providecommand{\subtitle}[1]{% add subtitle to \maketitle
  \apptocmd{\@title}{\par {\large #1 \par}}{}{}
}
\makeatother
\subtitle{dataset:}
\author{Magdalena A. Tkacz}
\date{}

\begin{document}
\maketitle

{
\setcounter{tocdepth}{3}
\tableofcontents
}
This document provides an example report structure based on the CRISP-DM
methodology.

Students may use, modify, or extend the template as appropriate for
their project.

The practical examples presented in this report are based on
\texttt{penguin\_mess}, a modified version of the
\texttt{palmerpenguins} dataset created for educational purposes.

The dataset was intentionally altered to include selected data quality
issues and should not be regarded as the original
\texttt{palmerpenguins} dataset.

\section{Assessment Summary}\label{assessment-summary}

Student name: Pen Guin

Intended grade:☐ 3.0 ☐ 3.5 ☑ 4.0 ☐ 4.5 ☐ 5.0

\section{Student Declaration}\label{student-declaration}

By submitting this report, I declare that:

\begin{itemize}
\tightlist
\item
  I understand all parts of the submitted work;
\item
  I am able to explain and justify the analytical decisions presented in
  the report;
\item
  the self-assessment table has been completed honestly;
\item
  any use of AI tools was limited to assistance and does not substitute
  for my own understanding of the analysis;
\item
  all external sources and tools used during the project have been
  properly acknowledged.
\end{itemize}

Submission of this report constitutes acceptance of this declaration.

\section{0. Executive Summary}\label{executive-summary}

\subsection{Report Summary}\label{report-summary}

This report presents a complete Data Mining workflow based on the
CRISP-DM methodology. The analysis uses the penguin\_mess dataset, a
modified version of the Palmer Penguins dataset prepared for educational
purposes.

The primary objective of the project is to predict penguin species using
morphological measurements and selected categorical attributes.

The analysis follows the CRISP-DM (Cross-Industry Standard Process for
Data Mining) methodology, progressing through business understanding,
data understanding, data preparation, exploratory analysis, modeling,
and evaluation.

Two classification models were developed and compared. The results
demonstrate the importance of proper data preparation and feature
selection when building predictive models.

\subsection{Success Criteria}\label{success-criteria}

The project will be considered successful if:

\begin{itemize}
\tightlist
\item
  data quality issues are identified and addressed;
\item
  meaningful exploratory analysis is performed;
\item
  at least two classification models are developed;
\item
  model performance is evaluated and compared;
\item
  conclusions are supported by analytical evidence.
\end{itemize}

\section{1. Business Understanding}\label{business-understanding}

The objective of this project is to develop a predictive model capable
of identifying penguin species based on available measurements and
characteristics.

Although the dataset is educational in nature, the problem reflects a
common Data Mining scenario in which historical observations are used to
classify new cases. The project also aims to demonstrate the complete
CRISP-DM process rather than maximize predictive performance.

The project aims to answer the following question:

\begin{quote}
Can penguin species be accurately predicted using the available
measurements and characteristics?
\end{quote}

\subsection{Project Scope}\label{project-scope}

The analysis includes: - data understanding and quality assessment; -
data preparation and cleaning; - exploratory data analysis (EDA); -
development of classification models; - model evaluation and comparison;
- interpretation of results.

The project does not aim to build a production-ready system. Instead, it
serves as a demonstration of the CRISP-DM process and fundamental Data
Mining techniques.

\subsection{Problem Statement}\label{problem-statement}

Given a set of penguin observations containing morphological
measurements and categorical characteristics, the objective is to
develop a predictive model capable of correctly identifying the species
of each penguin.

The project investigates which variables are most informative for
species classification and evaluates whether reliable predictions can be
obtained from the available data.

\subsection{Objective}\label{objective}

The objective of this project is to investigate whether penguin species
can be accurately predicted using available morphological measurements
and categorical attributes.

To achieve this goal, the dataset will be explored, cleaned, and
prepared for modeling. Multiple classification models will be developed
and compared in order to identify the most suitable approach for this
prediction task.

\subsection{Expected Outcome}\label{expected-outcome}

The expected outcome of this project is the development of a
classification model capable of predicting penguin species based on
available measurements and categorical attributes. It is expected that
the analysis will identify variables that are most informative for
species classification and reveal patterns that distinguish different
species.

Furthermore, the project is expected to demonstrate the impact of data
preparation and feature selection on model performance, as well as
provide a comparison of alternative modeling approaches.

\section{2. Data Understanding}\label{data-understanding}

The purpose of this phase is to become familiar with the dataset and
assess its suitability for the intended analysis. This includes
examining the structure of the data, identifying variable types,
evaluating data quality, and detecting potential issues that may affect
subsequent stages of the project.

Particular attention is given to missing values, inconsistencies,
unusual observations, and the distribution of individual variables.
Understanding these characteristics is essential before proceeding to
data preparation and model development.

\subsection{Dataset Description}\label{dataset-description}

The dataset was imported and inspected to understand its structure,
variable types, and overall characteristics.

An initial review was performed to identify the number of observations,
available variables, data types, and potential data quality issues that
may require attention during the preparation phase.

\begin{Shaded}
\begin{Highlighting}[]
\NormalTok{data }\OtherTok{\textless{}{-}} \FunctionTok{read\_csv}\NormalTok{(}
  \FunctionTok{here}\NormalTok{(}\StringTok{"data"}\NormalTok{, }\StringTok{"penguins\_mess.csv"}\NormalTok{)}
\NormalTok{)}
\end{Highlighting}
\end{Shaded}

\subsubsection{Basic Dataset
Information}\label{basic-dataset-information}

\begin{Shaded}
\begin{Highlighting}[]
\FunctionTok{dim}\NormalTok{(data)}
\end{Highlighting}
\end{Shaded}

\begin{verbatim}
## [1] 520  10
\end{verbatim}

So dataset has 520 observations and 10 variables.

Summary statistics

\begin{Shaded}
\begin{Highlighting}[]
\FunctionTok{summary}\NormalTok{(data)}
\end{Highlighting}
\end{Shaded}

\begin{verbatim}
##        id          species             island          bill_length_mm  
##  Min.   :  1.0   Length:520         Length:520         Min.   :  0.50  
##  1st Qu.:129.8   Class :character   Class :character   1st Qu.: 39.42  
##  Median :258.5   Mode  :character   Mode  :character   Median : 43.85  
##  Mean   :257.7                                         Mean   : 45.63  
##  3rd Qu.:385.2                                         3rd Qu.: 48.20  
##  Max.   :515.0                                         Max.   :999.00  
##                                                        NA's   :26      
##  bill_depth_mm   flipper_length_mm  body_mass_g        sex           
##  Min.   :12.60   Min.   :   1.0    Min.   : -500   Length:520        
##  1st Qu.:15.60   1st Qu.: 189.0    1st Qu.: 3563   Class :character  
##  Median :17.60   Median : 196.0    Median : 4043   Mode  :character  
##  Mean   :17.25   Mean   : 219.2    Mean   : 4332                     
##  3rd Qu.:18.80   3rd Qu.: 213.0    3rd Qu.: 4862                     
##  Max.   :21.70   Max.   :9999.0    Max.   :99000                     
##  NA's   :22      NA's   :20        NA's   :25                        
##       year         notes          
##  Min.   :1999   Length:520        
##  1st Qu.:2007   Class :character  
##  Median :2008   Mode  :character  
##  Mean   :2008                     
##  3rd Qu.:2009                     
##  Max.   :2099                     
## 
\end{verbatim}

\subsection{Dataset Structure}\label{dataset-structure}

\begin{Shaded}
\begin{Highlighting}[]
\FunctionTok{glimpse}\NormalTok{(data)}
\end{Highlighting}
\end{Shaded}

\begin{verbatim}
## Rows: 520
## Columns: 10
## $ id                <dbl> 4, 322, 105, 309, 69, 101, 334, 406, 58, 49, 467, 45~
## $ species           <chr> "Gentoo", "Adelie", "Gentoo", "Adelie", "Adelie", "A~
## $ island            <chr> "Biscoe", "Torgersen", "Biscoe", "Torgersen", "Torge~
## $ bill_length_mm    <dbl> 46.8, 35.9, 48.6, 37.9, 39.3, 41.9, 46.1, 41.8, 35.9~
## $ bill_depth_mm     <dbl> 15.4, 19.0, 14.7, 19.8, 19.6, 18.9, 14.6, 17.9, 17.5~
## $ flipper_length_mm <dbl> 217, 201, 217, 198, NA, 194, 208, 192, 193, 206, 192~
## $ body_mass_g       <dbl> 4706, 2886, 4644, 3514, 3514, 3080, 4874, 4222, 4254~
## $ sex               <chr> NA, "male", "female", "male", "female", "male", "mal~
## $ year              <dbl> 2009, 2007, 2009, 2007, 2009, 2007, 2007, 2008, 2009~
## $ notes             <chr> NA, NA, NA, NA, NA, NA, NA, NA, NA, NA, NA, NA, NA, ~
\end{verbatim}

\subsubsection{Variable Description}\label{variable-description}

The dataset contains a mixture of numerical and categorical variables.
Understanding the role and data type of each variable is important for
selecting appropriate preprocessing techniques, visualization methods,
and modeling approaches. The following table summarizes the variables
available in the dataset.

\begin{longtable}[]{@{}llr@{}}
\caption{Summary of variables and missing values.}\tabularnewline
\toprule\noalign{}
Variable & Data\_Type & Missing\_Values \\
\midrule\noalign{}
\endfirsthead
\toprule\noalign{}
Variable & Data\_Type & Missing\_Values \\
\midrule\noalign{}
\endhead
\bottomrule\noalign{}
\endlastfoot
id & numeric & 0 \\
species & character & 0 \\
island & character & 0 \\
bill\_length\_mm & numeric & 26 \\
bill\_depth\_mm & numeric & 22 \\
flipper\_length\_mm & numeric & 20 \\
body\_mass\_g & numeric & 25 \\
sex & character & 41 \\
year & numeric & 0 \\
notes & character & 484 \\
\end{longtable}

\subsection{Data Quality Assessment}\label{data-quality-assessment}

The dataset contains several missing values affecting both numerical and
categorical variables. Missing values are particularly common in the
\texttt{notes} attribute, where the majority of observations contain no
information.

In addition, the dataset includes an identifier variable (\texttt{id})
which is not expected to contribute to predictive performance and may
require removal before modeling.

The \texttt{notes} variable contains missing values in the majority of
observations and provides limited analytical value. Therefore, it was
excluded from further analysis.

\subsubsection{Missing Values}\label{missing-values}

Missing values were examined to assess the completeness of the dataset
and identify variables requiring additional preprocessing.

\begin{longtable}[]{@{}llrr@{}}
\caption{Missing values by variable.}\tabularnewline
\toprule\noalign{}
& Variable & Missing\_Values & Missing\_Percent \\
\midrule\noalign{}
\endfirsthead
\toprule\noalign{}
& Variable & Missing\_Values & Missing\_Percent \\
\midrule\noalign{}
\endhead
\bottomrule\noalign{}
\endlastfoot
notes & notes & 484 & 93.08 \\
sex & sex & 41 & 7.88 \\
bill\_length\_mm & bill\_length\_mm & 26 & 5.00 \\
body\_mass\_g & body\_mass\_g & 25 & 4.81 \\
bill\_depth\_mm & bill\_depth\_mm & 22 & 4.23 \\
flipper\_length\_mm & flipper\_length\_mm & 20 & 3.85 \\
id & id & 0 & 0.00 \\
species & species & 0 & 0.00 \\
island & island & 0 & 0.00 \\
year & year & 0 & 0.00 \\
\end{longtable}

The dataset contains missing values in several variables. The highest
proportion of missing data was observed in the \texttt{notes} attribute,
while selected numerical and categorical variables also contain
incomplete observations.

The impact of missing values will be assessed during the data
preparation phase, where appropriate handling strategies will be
applied.

\pandocbounded{\includegraphics[keepaspectratio]{example-report_files/figures/missing-values-plot-1.pdf}}
The analysis revealed missing values in both numerical and categorical
variables. The largest number of missing observations was found in the
\texttt{notes} attribute, whereas the remaining variables contain only a
relatively small proportion of incomplete records.

The identified missing values will be addressed during the data
preparation phase before model development.

\subsubsection{Duplicates}\label{duplicates}

\begin{Shaded}
\begin{Highlighting}[]
\FunctionTok{sum}\NormalTok{(}\FunctionTok{duplicated}\NormalTok{(data))}
\end{Highlighting}
\end{Shaded}

\begin{verbatim}
## [1] 5
\end{verbatim}

So, in more details:

\begin{Shaded}
\begin{Highlighting}[]
\NormalTok{data\_no\_id }\OtherTok{\textless{}{-}}\NormalTok{ dplyr}\SpecialCharTok{::}\FunctionTok{select}\NormalTok{(data, }\SpecialCharTok{{-}}\NormalTok{id)}

\FunctionTok{sum}\NormalTok{(}\FunctionTok{duplicated}\NormalTok{(data\_no\_id))}
\end{Highlighting}
\end{Shaded}

\begin{verbatim}
## [1] 19
\end{verbatim}

Several data quality issues were identified during the initial
assessment. Missing values are present in both numerical and categorical
variables, with the highest proportion observed in the \texttt{notes}
attribute.

In addition, a small number of duplicate observations were detected.
These records will be reviewed during the data preparation phase to
determine whether they should be retained or removed prior to model
development.

\subsubsection{Outliers}\label{outliers}

Potential outliers were investigated using graphical methods and summary
statistics. The objective was to identify observations that may
disproportionately influence subsequent analyses and predictive models.

\subsubsection{Outliers}\label{outliers-1}

\pandocbounded{\includegraphics[keepaspectratio]{example-report_files/figures/outliers-boxplots2-1.pdf}}
Several observations appear outside the whiskers of the boxplots. These
values were retained at this stage of the analysis and will be further
evaluated during data preparation.

\subsection{Initial Exploratory
Analysis}\label{initial-exploratory-analysis}

An initial exploratory analysis was conducted to examine the
distributions of individual variables and identify potential
relationships between them. The analysis focuses on understanding the
overall characteristics of the dataset and highlighting patterns that
may be relevant for subsequent modeling tasks.

\subsubsection{Distribution of Numerical
Variables}\label{distribution-of-numerical-variables}

Histograms were generated for all numerical variables to examine their
distributions, identify potential skewness, detect unusual observations,
and assess overall data quality. This step provides an initial
understanding of the range and shape of each variable before data
cleaning and modelling.

\begin{Shaded}
\begin{Highlighting}[]
\NormalTok{data }\SpecialCharTok{|\textgreater{}}
  \FunctionTok{select}\NormalTok{(}\FunctionTok{where}\NormalTok{(is.numeric)) }\SpecialCharTok{|\textgreater{}}
  \FunctionTok{pivot\_longer}\NormalTok{(}\FunctionTok{everything}\NormalTok{()) }\SpecialCharTok{|\textgreater{}}
  \FunctionTok{ggplot}\NormalTok{(}\FunctionTok{aes}\NormalTok{(value)) }\SpecialCharTok{+}
  \FunctionTok{geom\_histogram}\NormalTok{(}\AttributeTok{bins =} \DecValTok{20}\NormalTok{) }\SpecialCharTok{+}
  \FunctionTok{facet\_wrap}\NormalTok{(}\SpecialCharTok{\textasciitilde{}}\NormalTok{name, }\AttributeTok{scales =} \StringTok{"free"}\NormalTok{)}\SpecialCharTok{+}
  \FunctionTok{theme\_minimal}\NormalTok{()}
\end{Highlighting}
\end{Shaded}

\pandocbounded{\includegraphics[keepaspectratio]{example-report_files/figures/unnamed-chunk-7-1.pdf}}
The distributions of numerical variables exhibit several irregularities,
including missing values and extreme observations. These characteristics
may affect subsequent analyses and limit the interpretability of
visualizations and predictive models. The observed patterns indicate
that data preparation will be required before proceeding to model
development.

\subsubsection{Distribution of Categorical
Variables}\label{distribution-of-categorical-variables}

Distribution of Categorical Variables The distribution of categorical
variables was examined to understand the composition of the dataset and
identify potential imbalances between categories. Bar charts were used
to visualize the frequency of observations within each category.

\begin{Shaded}
\begin{Highlighting}[]
\NormalTok{data }\SpecialCharTok{|\textgreater{}}
  \FunctionTok{select}\NormalTok{(}\FunctionTok{where}\NormalTok{(is.character)) }\SpecialCharTok{|\textgreater{}}
  \FunctionTok{pivot\_longer}\NormalTok{(}\FunctionTok{everything}\NormalTok{()) }\SpecialCharTok{|\textgreater{}}
  \FunctionTok{ggplot}\NormalTok{(}\FunctionTok{aes}\NormalTok{(value)) }\SpecialCharTok{+}
  \FunctionTok{geom\_bar}\NormalTok{() }\SpecialCharTok{+}
  \FunctionTok{facet\_wrap}\NormalTok{(}\SpecialCharTok{\textasciitilde{}}\NormalTok{name, }\AttributeTok{scales =} \StringTok{"free"}\NormalTok{)}
\end{Highlighting}
\end{Shaded}

\pandocbounded{\includegraphics[keepaspectratio]{example-report_files/figures/unnamed-chunk-8-1.pdf}}
The examined categorical variables show varying frequencies across
categories. These distributions contribute to a better understanding of
the dataset and may be relevant during model development and evaluation.

\subsubsection{Relationships and Correlation Between
Variables}\label{relationships-and-correlation-between-variables}

Relationships between numerical variables were explored using a
scatterplot and scatterplot matrix. This visualization allows pairwise
comparison of a certain or all numerical attributes and helps identify
potential correlations, clusters, trends, and unusual observations that
may influence subsequent modelling steps.

\begin{Shaded}
\begin{Highlighting}[]
\FunctionTok{ggplot}\NormalTok{(}
\NormalTok{  data,}
  \FunctionTok{aes}\NormalTok{(}
\NormalTok{    bill\_length\_mm,}
\NormalTok{    flipper\_length\_mm,}
    \AttributeTok{color =}\NormalTok{ species}
\NormalTok{  )}
\NormalTok{) }\SpecialCharTok{+}
  \FunctionTok{geom\_point}\NormalTok{()}
\end{Highlighting}
\end{Shaded}

\pandocbounded{\includegraphics[keepaspectratio]{example-report_files/figures/unnamed-chunk-9-1.pdf}}
The initial scatterplots were strongly influenced by extreme
observations, making it difficult to visually assess relationships
between variables. This suggests the presence of potential outliers that
may require further investigation before meaningful interpretation can
be performed.

\begin{Shaded}
\begin{Highlighting}[]
\FunctionTok{library}\NormalTok{(GGally)}

\NormalTok{data }\SpecialCharTok{|\textgreater{}}
  \FunctionTok{select}\NormalTok{(}\FunctionTok{where}\NormalTok{(is.numeric)) }\SpecialCharTok{|\textgreater{}}
\NormalTok{  GGally}\SpecialCharTok{::}\FunctionTok{ggpairs}\NormalTok{(}
    \AttributeTok{progress =} \ConstantTok{FALSE}\NormalTok{,}
    \AttributeTok{title =} \StringTok{"Scatterplot Matrix of Numerical Variables"}
\NormalTok{  )}
\end{Highlighting}
\end{Shaded}

\pandocbounded{\includegraphics[keepaspectratio]{example-report_files/figures/unnamed-chunk-10-1.pdf}}

The visualization suggests that some numerical variables may be related.
A more detailed investigation is required to determine whether these
relationships are meaningful and relevant for predictive modelling.

The scatterplot matrix serves as an initial assessment of relationships
between numerical variables. Visual inspection suggests that some
variables may be correlated, although the strength and significance of
these relationships cannot be reliably determined from the plots alone.
Therefore, a dedicated correlation analysis is recommended to support
feature selection and improve understanding of the dataset structure.

\subsubsection{Correlation Analysis}\label{correlation-analysis}

Correlation analysis was performed to quantify relationships between
numerical variables. Understanding these relationships can support
feature selection, help identify redundant information, and provide
additional insight into the structure of the dataset.

\begin{Shaded}
\begin{Highlighting}[]
\FunctionTok{library}\NormalTok{(corrplot)}

\NormalTok{cor\_matrix }\OtherTok{\textless{}{-}}\NormalTok{ data }\SpecialCharTok{|\textgreater{}}
  \FunctionTok{select}\NormalTok{(}\FunctionTok{where}\NormalTok{(is.numeric)) }\SpecialCharTok{|\textgreater{}}
  \FunctionTok{cor}\NormalTok{(}\AttributeTok{use =} \StringTok{"pairwise.complete.obs"}\NormalTok{)}

\FunctionTok{corrplot}\NormalTok{(}
\NormalTok{  cor\_matrix,}
  \AttributeTok{method =} \StringTok{"color"}\NormalTok{,}
  \AttributeTok{type =} \StringTok{"upper"}\NormalTok{,}
  \AttributeTok{addCoef.col =} \StringTok{"black"}\NormalTok{,}
  \AttributeTok{tl.col =} \StringTok{"black"}\NormalTok{,}
  \AttributeTok{tl.srt =} \DecValTok{45}
\NormalTok{)}
\end{Highlighting}
\end{Shaded}

\pandocbounded{\includegraphics[keepaspectratio]{example-report_files/figures/unnamed-chunk-11-1.pdf}}

The results indicate varying levels of correlation among the numerical
variables. While some relationships appear stronger than others, the
practical significance of these associations should be evaluated within
the context of the modelling task. Correlation analysis serves as a
useful exploratory tool but should not be used as the sole criterion for
feature selection.

The correlation matrix provides a quantitative summary of linear
relationships between numerical variables. Some variables exhibit
stronger associations than others, suggesting that they may contain
related information. Correlation analysis can help identify potentially
redundant features and support feature selection during model
development. However, correlation measures only linear relationships and
should always be interpreted alongside exploratory visualisations and
domain knowledge.

ℹ️️ \textbf{Correlation does not imply dependence}

Correlation measures only linear relationships between variables. A low
correlation coefficient does not necessarily indicate the absence of a
relationship.

Variables may exhibit non-linear dependencies that are not captured by
correlation analysis. A classic example is the \emph{Anscombe Quartet},
where several datasets share nearly identical summary statistics and
correlation coefficients while displaying very different patterns when
visualized.

For this reason, correlation analysis should always be complemented by
exploratory visualisations and domain knowledge.

\subsection{Initial Observations}\label{initial-observations}

The dataset contains both numerical and categorical variables. Initial
inspection suggests that some data cleaning will be necessary due to
missing values and potential inconsistencies. Further exploratory
analysis is required to better understand the structure and quality of
the data.

\section{3. Data Preparation}\label{data-preparation}

The exploratory analysis revealed several data quality issues that may
affect the performance of predictive models. Therefore, a data cleaning
process was performed to improve consistency, reduce noise, and prepare
the dataset for further analysis and modelling.

\subsection{Missing Values}\label{missing-values-1}

Missing values were identified in several variables. The notes column
contained mostly missing observations and limited analytical value;
therefore, it was removed from further analysis. Remaining missing
values will be handled using appropriate preprocessing techniques during
model development.

Removing \texttt{notes}:

\begin{Shaded}
\begin{Highlighting}[]
\NormalTok{data }\OtherTok{\textless{}{-}}\NormalTok{ data }\SpecialCharTok{|\textgreater{}}
  \FunctionTok{select}\NormalTok{(}\SpecialCharTok{{-}}\NormalTok{notes)}
\end{Highlighting}
\end{Shaded}

Drop na in \texttt{species}:

\begin{Shaded}
\begin{Highlighting}[]
\NormalTok{data }\OtherTok{\textless{}{-}}\NormalTok{ data }\SpecialCharTok{|\textgreater{}}
  \FunctionTok{drop\_na}\NormalTok{(species)}
\end{Highlighting}
\end{Shaded}

\subsection{Duplicates}\label{duplicates-1}

Duplicate observations were identified and removed to prevent potential
bias during model training and evaluation.

\begin{Shaded}
\begin{Highlighting}[]
\NormalTok{data }\OtherTok{\textless{}{-}}\NormalTok{ data }\SpecialCharTok{|\textgreater{}}
  \FunctionTok{distinct}\NormalTok{()}
\end{Highlighting}
\end{Shaded}

\subsection{Data Cleaning}\label{data-cleaning}

\subsubsection{Outliers}\label{outliers-2}

\begin{Shaded}
\begin{Highlighting}[]
\FunctionTok{summary}\NormalTok{(data)}
\end{Highlighting}
\end{Shaded}

\begin{verbatim}
##        id          species             island          bill_length_mm  
##  Min.   :  1.0   Length:515         Length:515         Min.   :  0.50  
##  1st Qu.:129.5   Class :character   Class :character   1st Qu.: 39.40  
##  Median :258.0   Mode  :character   Mode  :character   Median : 43.80  
##  Mean   :258.0                                         Mean   : 45.62  
##  3rd Qu.:386.5                                         3rd Qu.: 48.20  
##  Max.   :515.0                                         Max.   :999.00  
##                                                        NA's   :25      
##  bill_depth_mm   flipper_length_mm  body_mass_g        sex           
##  Min.   :12.60   Min.   :   1.0    Min.   : -500   Length:515        
##  1st Qu.:15.60   1st Qu.: 189.0    1st Qu.: 3547   Class :character  
##  Median :17.60   Median : 196.0    Median : 4041   Mode  :character  
##  Mean   :17.25   Mean   : 219.3    Mean   : 4328                     
##  3rd Qu.:18.80   3rd Qu.: 213.0    3rd Qu.: 4852                     
##  Max.   :21.70   Max.   :9999.0    Max.   :99000                     
##  NA's   :21      NA's   :20        NA's   :25                        
##       year     
##  Min.   :1999  
##  1st Qu.:2007  
##  Median :2008  
##  Mean   :2008  
##  3rd Qu.:2009  
##  Max.   :2099  
## 
\end{verbatim}

Extreme values outside plausible biological ranges were identified
during exploratory analysis. These observations were considered data
quality issues and were removed prior to modelling.

\begin{Shaded}
\begin{Highlighting}[]
\NormalTok{data }\OtherTok{\textless{}{-}}\NormalTok{ data }\SpecialCharTok{|\textgreater{}}
  \FunctionTok{filter}\NormalTok{(}
\NormalTok{    bill\_length\_mm }\SpecialCharTok{\textless{}} \DecValTok{100}\NormalTok{,}
\NormalTok{    flipper\_length\_mm }\SpecialCharTok{\textless{}} \DecValTok{400}\NormalTok{,}
\NormalTok{    body\_mass\_g }\SpecialCharTok{\textless{}} \DecValTok{10000}\NormalTok{,}
\NormalTok{    year }\SpecialCharTok{\textless{}} \DecValTok{2100}
\NormalTok{  )}
\end{Highlighting}
\end{Shaded}

\subsection{Transformations}\label{transformations}

\subsubsection{Standardizing Categorical
Variables}\label{standardizing-categorical-variables}

\begin{Shaded}
\begin{Highlighting}[]
\NormalTok{data }\OtherTok{\textless{}{-}}\NormalTok{ data }\SpecialCharTok{|\textgreater{}}
  \FunctionTok{mutate}\NormalTok{(}
    \AttributeTok{species =} \FunctionTok{str\_to\_title}\NormalTok{(species),}
    \AttributeTok{island =} \FunctionTok{str\_to\_title}\NormalTok{(island),}
    \AttributeTok{sex =} \FunctionTok{str\_to\_lower}\NormalTok{(sex)}
\NormalTok{  )}
\end{Highlighting}
\end{Shaded}

And

\begin{Shaded}
\begin{Highlighting}[]
\NormalTok{data }\OtherTok{\textless{}{-}}\NormalTok{ data }\SpecialCharTok{|\textgreater{}}
  \FunctionTok{mutate}\NormalTok{(}
    \AttributeTok{species =} \FunctionTok{recode}\NormalTok{(}
\NormalTok{      species,}
      \StringTok{"Gent00"} \OtherTok{=} \StringTok{"Gentoo"}
\NormalTok{    ),}
    \AttributeTok{sex =} \FunctionTok{recode}\NormalTok{(}
\NormalTok{      sex,}
      \StringTok{"m"} \OtherTok{=} \StringTok{"male"}\NormalTok{,}
      \StringTok{"f"} \OtherTok{=} \StringTok{"female"}
\NormalTok{    )}
\NormalTok{  )}
\end{Highlighting}
\end{Shaded}

Several categorical variables contained inconsistent labels resulting
from differences in spelling, capitalisation, and data entry errors.
Categories were standardised to ensure consistent representation across
the dataset.

\subsection{Feature Selection /
Engineering}\label{feature-selection-engineering}

Several modifications were applied during the data preparation stage.

The \texttt{notes} variable was removed because it contained a large
proportion of missing values and provided limited analytical value.
Duplicate records were removed to improve data consistency.

Categorical variables were standardized to eliminate inconsistencies
caused by differences in spelling, capitalization, and data entry
errors. Extreme observations identified as likely data quality issues
were removed from the dataset.

No new variables were created at this stage. Feature engineering will be
limited to preprocessing steps required for model development, such as
encoding categorical variables and handling remaining missing values.

Potential additional feature engineering techniques, such as creating
derived variables or applying transformations, may be explored depending
on the modelling approach and project objectives.

\subsection{Check}\label{check}

\begin{Shaded}
\begin{Highlighting}[]
\FunctionTok{glimpse}\NormalTok{(data)}
\end{Highlighting}
\end{Shaded}

\begin{verbatim}
## Rows: 443
## Columns: 9
## $ id                <dbl> 4, 322, 105, 309, 101, 334, 406, 58, 49, 467, 453, 4~
## $ species           <chr> "Gentoo", "Adelie", "Gentoo", "Adelie", "Adelie", "G~
## $ island            <chr> "Biscoe", "Torgersen", "Biscoe", "Torgersen", "Dream~
## $ bill_length_mm    <dbl> 46.8, 35.9, 48.6, 37.9, 41.9, 46.1, 41.8, 35.9, 44.3~
## $ bill_depth_mm     <dbl> 15.4, 19.0, 14.7, 19.8, 18.9, 14.6, 17.9, 17.5, 16.2~
## $ flipper_length_mm <dbl> 217, 201, 217, 198, 194, 208, 192, 193, 206, 192, 17~
## $ body_mass_g       <dbl> 4706, 2886, 4644, 3514, 3080, 4874, 4222, 4254, 4764~
## $ sex               <chr> NA, "male", "female", "male", "male", "male", "male"~
## $ year              <dbl> 2009, 2007, 2009, 2007, 2007, 2007, 2008, 2009, 2009~
\end{verbatim}

\begin{Shaded}
\begin{Highlighting}[]
\FunctionTok{summary}\NormalTok{(data)}
\end{Highlighting}
\end{Shaded}

\begin{verbatim}
##        id          species             island          bill_length_mm 
##  Min.   :  2.0   Length:443         Length:443         Min.   :31.20  
##  1st Qu.:132.5   Class :character   Class :character   1st Qu.:39.50  
##  Median :258.0   Mode  :character   Mode  :character   Median :43.60  
##  Mean   :259.5                                         Mean   :43.69  
##  3rd Qu.:386.5                                         3rd Qu.:48.10  
##  Max.   :514.0                                         Max.   :58.30  
##                                                                       
##  bill_depth_mm   flipper_length_mm  body_mass_g       sex           
##  Min.   :12.60   Min.   :  1.0     Min.   :-500   Length:443        
##  1st Qu.:15.70   1st Qu.:189.0     1st Qu.:3563   Class :character  
##  Median :17.60   Median :196.0     Median :4042   Mode  :character  
##  Mean   :17.25   Mean   :199.6     Mean   :4130                     
##  3rd Qu.:18.70   3rd Qu.:213.0     3rd Qu.:4822                     
##  Max.   :21.70   Max.   :232.0     Max.   :6843                     
##  NA's   :20                                                         
##       year     
##  Min.   :2007  
##  1st Qu.:2007  
##  Median :2008  
##  Mean   :2008  
##  3rd Qu.:2009  
##  Max.   :2099  
## 
\end{verbatim}

Category labels were standardized; however, some missing values remained
in the sex variable. These observations were retained at this stage and
will be handled during the modelling phase using an appropriate
preprocessing strategy.

\subsubsection{Missing Values}\label{missing-values-2}

Missing values were examined to assess the completeness of the dataset
and identify variables requiring additional preprocessing.

\begin{longtable}[]{@{}llrr@{}}
\caption{Missing values by variable.}\tabularnewline
\toprule\noalign{}
& Variable & Missing\_Values & Missing\_Percent \\
\midrule\noalign{}
\endfirsthead
\toprule\noalign{}
& Variable & Missing\_Values & Missing\_Percent \\
\midrule\noalign{}
\endhead
\bottomrule\noalign{}
\endlastfoot
sex & sex & 36 & 8.13 \\
bill\_depth\_mm & bill\_depth\_mm & 20 & 4.51 \\
id & id & 0 & 0.00 \\
species & species & 0 & 0.00 \\
island & island & 0 & 0.00 \\
bill\_length\_mm & bill\_length\_mm & 0 & 0.00 \\
flipper\_length\_mm & flipper\_length\_mm & 0 & 0.00 \\
body\_mass\_g & body\_mass\_g & 0 & 0.00 \\
year & year & 0 & 0.00 \\
\end{longtable}

The dataset still contains missing values in two variables. The impact
of missing values will be assessed during the data preparation before
modeling phase, where appropriate handling strategies will be applied.

\pandocbounded{\includegraphics[keepaspectratio]{example-report_files/figures/missing-values2-plot-1.pdf}}

\subsubsection{Duplicates}\label{duplicates-2}

\begin{Shaded}
\begin{Highlighting}[]
\FunctionTok{sum}\NormalTok{(}\FunctionTok{duplicated}\NormalTok{(data))}
\end{Highlighting}
\end{Shaded}

\begin{verbatim}
## [1] 0
\end{verbatim}

\subsubsection{Outliers}\label{outliers-3}

\pandocbounded{\includegraphics[keepaspectratio]{example-report_files/figures/outliers-boxplots-1.pdf}}

\section{4. Exploratory Data Analysis
(EDA)}\label{exploratory-data-analysis-eda}

\subsection{Distribution Analysis}\label{distribution-analysis}

\begin{Shaded}
\begin{Highlighting}[]
\NormalTok{data }\SpecialCharTok{|\textgreater{}}
  \FunctionTok{select}\NormalTok{(}\FunctionTok{where}\NormalTok{(is.numeric)) }\SpecialCharTok{|\textgreater{}}
  \FunctionTok{pivot\_longer}\NormalTok{(}\FunctionTok{everything}\NormalTok{()) }\SpecialCharTok{|\textgreater{}}
  \FunctionTok{ggplot}\NormalTok{(}\FunctionTok{aes}\NormalTok{(value)) }\SpecialCharTok{+}
  \FunctionTok{geom\_histogram}\NormalTok{(}\AttributeTok{bins =} \DecValTok{20}\NormalTok{) }\SpecialCharTok{+}
  \FunctionTok{facet\_wrap}\NormalTok{(}\SpecialCharTok{\textasciitilde{}}\NormalTok{name, }\AttributeTok{scales =} \StringTok{"free"}\NormalTok{)}
\end{Highlighting}
\end{Shaded}

\pandocbounded{\includegraphics[keepaspectratio]{example-report_files/figures/unnamed-chunk-21-1.pdf}}

\subsubsection{Interpretation}\label{interpretation}

\begin{quote}
The distributions provide an overview of the numerical variables and
their ranges. Several variables exhibit unusual observations that may
require further investigation. Understanding these distributions is
important for selecting appropriate preprocessing and modelling
techniques.
\end{quote}

\subsection{Relationship Analysis}\label{relationship-analysis}

\textbf{Introduction}

\begin{quote}
Relationships between numerical variables were explored using a
scatterplot matrix. This visualization provides an overview of pairwise
interactions and may reveal patterns worth investigating in later stages
of the analysis.
\end{quote}

\begin{Shaded}
\begin{Highlighting}[]
\FunctionTok{library}\NormalTok{(GGally)}

\NormalTok{data }\SpecialCharTok{|\textgreater{}}
  \FunctionTok{select}\NormalTok{(}\FunctionTok{where}\NormalTok{(is.numeric)) }\SpecialCharTok{|\textgreater{}}
  \FunctionTok{ggpairs}\NormalTok{(}\AttributeTok{progress =} \ConstantTok{FALSE}\NormalTok{)}
\end{Highlighting}
\end{Shaded}

\pandocbounded{\includegraphics[keepaspectratio]{example-report_files/figures/unnamed-chunk-22-1.pdf}}

\subsubsection{Interpretation}\label{interpretation-1}

\begin{quote}
The scatterplot matrix provides an initial overview of potential
relationships among variables. Several patterns can be observed and may
warrant further investigation during subsequent exploratory analysis and
feature selection stages.
\end{quote}

\subsection{Correlation Analysis}\label{correlation-analysis-1}

\textbf{Introduction}

\begin{quote}
Correlation analysis was performed to quantify linear relationships
between numerical variables. Understanding these relationships can
support feature selection and provide additional insight into the
structure of the dataset.
\end{quote}

\begin{Shaded}
\begin{Highlighting}[]
\FunctionTok{library}\NormalTok{(corrplot)}

\NormalTok{cor\_matrix }\OtherTok{\textless{}{-}}\NormalTok{ data }\SpecialCharTok{|\textgreater{}}
  \FunctionTok{select}\NormalTok{(}\FunctionTok{where}\NormalTok{(is.numeric)) }\SpecialCharTok{|\textgreater{}}
  \FunctionTok{cor}\NormalTok{(}\AttributeTok{use =} \StringTok{"pairwise.complete.obs"}\NormalTok{)}

\FunctionTok{corrplot}\NormalTok{(}
\NormalTok{  cor\_matrix,}
  \AttributeTok{method =} \StringTok{"color"}\NormalTok{,}
  \AttributeTok{type =} \StringTok{"upper"}\NormalTok{,}
  \AttributeTok{addCoef.col =} \StringTok{"black"}\NormalTok{,}
  \AttributeTok{tl.col =} \StringTok{"black"}\NormalTok{,}
  \AttributeTok{tl.srt =} \DecValTok{45}
\NormalTok{)}
\end{Highlighting}
\end{Shaded}

\pandocbounded{\includegraphics[keepaspectratio]{example-report_files/figures/unnamed-chunk-23-1.pdf}}

\subsubsection{Interpretation}\label{interpretation-2}

\begin{quote}
The results indicate varying levels of correlation among the numerical
variables. While some relationships appear stronger than others, the
practical significance of these associations should be evaluated within
the context of the modelling task. Correlation analysis serves as a
useful exploratory tool but should not be used as the sole criterion for
feature selection.
\end{quote}

\begin{Shaded}
\begin{Highlighting}[]
\DataTypeTok{\textless{}}\KeywordTok{div}\OtherTok{ style}\OperatorTok{=}\StringTok{"}
\StringTok{background{-}color:\#e8f1f8;}
\StringTok{border{-}left:8px solid \#5c7fa3;}
\StringTok{color:\#111111;}
\StringTok{padding:12px 16px;}
\StringTok{margin:12px 0;}
\StringTok{"}\DataTypeTok{\textgreater{}}

\NormalTok{💡 **Correlation does not imply dependence**}

\NormalTok{Correlation measures only linear relationships between variables. Variables with low correlation may still be related through non{-}linear patterns. The *Anscombe Quartet* is a well{-}known example demonstrating why visual inspection should accompany statistical summaries.}

\DataTypeTok{\textless{}/}\KeywordTok{div}\DataTypeTok{\textgreater{}}
\end{Highlighting}
\end{Shaded}

\subsection{Outlier Analysis}\label{outlier-analysis}

\textbf{Introduction}

\begin{quote}
Outlier analysis was performed to identify observations that differ
substantially from the majority of the data. Such observations may
represent valid extreme cases or potential data quality issues.
\end{quote}

\begin{Shaded}
\begin{Highlighting}[]
\NormalTok{data }\SpecialCharTok{|\textgreater{}}
  \FunctionTok{select}\NormalTok{(}\FunctionTok{where}\NormalTok{(is.numeric)) }\SpecialCharTok{|\textgreater{}}
  \FunctionTok{pivot\_longer}\NormalTok{(}\FunctionTok{everything}\NormalTok{()) }\SpecialCharTok{|\textgreater{}}
  \FunctionTok{ggplot}\NormalTok{(}\FunctionTok{aes}\NormalTok{(}\AttributeTok{y =}\NormalTok{ value, }\AttributeTok{x =}\NormalTok{ name)) }\SpecialCharTok{+}
  \FunctionTok{geom\_boxplot}\NormalTok{() }\SpecialCharTok{+}
  \FunctionTok{coord\_flip}\NormalTok{()}
\end{Highlighting}
\end{Shaded}

\pandocbounded{\includegraphics[keepaspectratio]{example-report_files/figures/unnamed-chunk-24-1.pdf}}

\subsubsection{Interpretation}\label{interpretation-3}

\begin{quote}
Several observations fall outside the typical range of values observed
in the dataset. These observations should be reviewed carefully, as they
may influence statistical summaries and predictive model performance.
\end{quote}

\subsection{Exploring Body Mass
Differences}\label{exploring-body-mass-differences}

Exploratory data analysis is not limited to relationships between
numerical variables. Relationships between categorical and numerical
variables can also provide valuable insights. The following example
examines differences in body mass across sex categories.

\begin{Shaded}
\begin{Highlighting}[]
\NormalTok{data }\SpecialCharTok{|\textgreater{}}
  \FunctionTok{filter}\NormalTok{(}\SpecialCharTok{!}\FunctionTok{is.na}\NormalTok{(sex)) }\SpecialCharTok{|\textgreater{}}
  \FunctionTok{ggplot}\NormalTok{(}
    \FunctionTok{aes}\NormalTok{(}
      \AttributeTok{x =}\NormalTok{ sex,}
      \AttributeTok{y =}\NormalTok{ body\_mass\_g,}
      \AttributeTok{fill =}\NormalTok{ sex}
\NormalTok{    )}
\NormalTok{  ) }\SpecialCharTok{+}
  \FunctionTok{geom\_violin}\NormalTok{(}\AttributeTok{trim =} \ConstantTok{FALSE}\NormalTok{) }\SpecialCharTok{+}
  \FunctionTok{geom\_boxplot}\NormalTok{(}
    \AttributeTok{width =} \FloatTok{0.15}\NormalTok{,}
    \AttributeTok{fill =} \StringTok{"white"}
\NormalTok{  ) }\SpecialCharTok{+}
  \FunctionTok{theme\_minimal}\NormalTok{()}
\end{Highlighting}
\end{Shaded}

\pandocbounded{\includegraphics[keepaspectratio]{example-report_files/figures/unnamed-chunk-25-1.pdf}}

\begin{Shaded}
\begin{Highlighting}[]
\NormalTok{data }\SpecialCharTok{|\textgreater{}}
  \FunctionTok{filter}\NormalTok{(}\SpecialCharTok{!}\FunctionTok{is.na}\NormalTok{(sex)) }\SpecialCharTok{|\textgreater{}}
  \FunctionTok{group\_by}\NormalTok{(sex) }\SpecialCharTok{|\textgreater{}}
  \FunctionTok{summarise}\NormalTok{(}
    \AttributeTok{Mean\_Body\_Mass =} \FunctionTok{mean}\NormalTok{(}
\NormalTok{      body\_mass\_g,}
      \AttributeTok{na.rm =} \ConstantTok{TRUE}
\NormalTok{    ),}
    \AttributeTok{SD =} \FunctionTok{sd}\NormalTok{(}
\NormalTok{      body\_mass\_g,}
      \AttributeTok{na.rm =} \ConstantTok{TRUE}
\NormalTok{    )}
\NormalTok{  )}
\end{Highlighting}
\end{Shaded}

\begin{verbatim}
## # A tibble: 2 x 3
##   sex    Mean_Body_Mass    SD
##   <chr>           <dbl> <dbl>
## 1 female          4083. 1089.
## 2 male            4142.  891.
\end{verbatim}

Exploratory analysis suggests that body mass differs between male and
female penguins. Although the distributions overlap, male penguins tend
to have higher body mass on average. This observation indicates that sex
may provide useful information for predictive modelling.

\subsection{Summary of EDA}\label{summary-of-eda}

\begin{quote}
The exploratory analysis provided a better understanding of the dataset
structure and quality. Several data quality issues were identified,
including missing values, inconsistent category labels, duplicate
records, and unusual observations. The analysis also revealed potential
relationships between variables that may be useful for predictive
modelling. These findings support the data preparation decisions
described in the previous section and provide a foundation for model
development.
\end{quote}

\section{5. Modeling}\label{modeling}

\subsection{Data Preparation Strategy}\label{data-preparation-strategy}

The exploratory analysis revealed several data quality issues, including
missing values, duplicate records, inconsistent categorical labels, and
extreme observations. Appropriate preprocessing steps were therefore
applied to improve data consistency and prepare the dataset for machine
learning models.

\begin{Shaded}
\begin{Highlighting}[]
\FunctionTok{library}\NormalTok{(tidyverse)}
\FunctionTok{library}\NormalTok{(tidymodels)}
\end{Highlighting}
\end{Shaded}

\subsection{Selected Features}\label{selected-features}

The target variable for this project is \texttt{species}.

The following variables were retained as potential predictors:

\begin{itemize}
\tightlist
\item
  bill\_length\_mm
\item
  bill\_depth\_mm
\item
  flipper\_length\_mm
\item
  body\_mass\_g
\item
  island
\item
  sex
\item
  year
\end{itemize}

The \texttt{notes} variable was excluded from further analysis due to
its large number of missing values and limited analytical relevance.

\begin{Shaded}
\begin{Highlighting}[]
\NormalTok{model\_data }\OtherTok{\textless{}{-}}\NormalTok{ data }\SpecialCharTok{|\textgreater{}}
  \FunctionTok{select}\NormalTok{(}
\NormalTok{    species,}
\NormalTok{    bill\_length\_mm,}
\NormalTok{    bill\_depth\_mm,}
\NormalTok{    flipper\_length\_mm,}
\NormalTok{    body\_mass\_g,}
\NormalTok{    island,}
\NormalTok{    sex,}
\NormalTok{    year}
\NormalTok{  )}
\end{Highlighting}
\end{Shaded}

\subsection{Data Cleaning}\label{data-cleaning-1}

Duplicate observations were removed and categorical variables were
standardized to ensure consistent representation across the dataset.
Obvious data quality issues and implausible observations identified
during exploratory analysis were addressed prior to modelling.

\begin{Shaded}
\begin{Highlighting}[]
\NormalTok{model\_data }\OtherTok{\textless{}{-}}\NormalTok{ model\_data }\SpecialCharTok{|\textgreater{}}
  \FunctionTok{distinct}\NormalTok{() }\SpecialCharTok{|\textgreater{}}
  \FunctionTok{mutate}\NormalTok{(}
    \AttributeTok{species =} \FunctionTok{str\_to\_title}\NormalTok{(species),}
    \AttributeTok{island =} \FunctionTok{str\_to\_title}\NormalTok{(island),}
    \AttributeTok{sex =} \FunctionTok{str\_to\_lower}\NormalTok{(sex),}
    \AttributeTok{sex =} \FunctionTok{recode}\NormalTok{(}
\NormalTok{      sex,}
      \StringTok{"m"} \OtherTok{=} \StringTok{"male"}\NormalTok{,}
      \StringTok{"f"} \OtherTok{=} \StringTok{"female"}
\NormalTok{    ),}
    \AttributeTok{species =} \FunctionTok{recode}\NormalTok{(}
\NormalTok{      species,}
      \StringTok{"Gent00"} \OtherTok{=} \StringTok{"Gentoo"}
\NormalTok{    )}
\NormalTok{  ) }\SpecialCharTok{|\textgreater{}}
  \FunctionTok{filter}\NormalTok{(}
\NormalTok{    bill\_length\_mm }\SpecialCharTok{\textless{}} \DecValTok{100} \SpecialCharTok{|} \FunctionTok{is.na}\NormalTok{(bill\_length\_mm),}
\NormalTok{    flipper\_length\_mm }\SpecialCharTok{\textless{}} \DecValTok{400} \SpecialCharTok{|} \FunctionTok{is.na}\NormalTok{(flipper\_length\_mm),}
\NormalTok{    body\_mass\_g }\SpecialCharTok{\textless{}} \DecValTok{10000} \SpecialCharTok{|} \FunctionTok{is.na}\NormalTok{(body\_mass\_g),}
\NormalTok{    year }\SpecialCharTok{\textless{}} \DecValTok{2100} \SpecialCharTok{|} \FunctionTok{is.na}\NormalTok{(year)}
\NormalTok{  )}
\end{Highlighting}
\end{Shaded}

\subsection{Handling Missing Values}\label{handling-missing-values}

Missing values remained in selected variables after the cleaning
process. These observations were retained and handled through
model-specific preprocessing techniques during model training.

Three possible strategies may be considered for the \texttt{sex}
variable:

\begin{itemize}
\tightlist
\item
  removing observations with missing values,
\item
  replacing missing values with a new category,
\item
  imputing missing values during model training.
\end{itemize}

In this report, the third approach is preferred because it keeps more
observations in the dataset and allows preprocessing to be included
directly in the modelling workflow.

\subsubsection{\texorpdfstring{Option 1: Removing Records with Missing
\texttt{sex}}{Option 1: Removing Records with Missing sex}}\label{option-1-removing-records-with-missing-sex}

\begin{Shaded}
\begin{Highlighting}[]
\NormalTok{model\_data\_option\_1 }\OtherTok{\textless{}{-}}\NormalTok{ model\_data }\SpecialCharTok{|\textgreater{}}
  \FunctionTok{drop\_na}\NormalTok{(sex)}
\end{Highlighting}
\end{Shaded}

This approach is simple, but it reduces the size of the dataset. It may
be acceptable when the number of missing observations is small, but it
can also remove potentially useful information.

\subsubsection{\texorpdfstring{Option 2: Creating an \texttt{unknown}
Category}{Option 2: Creating an unknown Category}}\label{option-2-creating-an-unknown-category}

\begin{Shaded}
\begin{Highlighting}[]
\NormalTok{model\_data\_option\_2 }\OtherTok{\textless{}{-}}\NormalTok{ model\_data }\SpecialCharTok{|\textgreater{}}
  \FunctionTok{mutate}\NormalTok{(}
    \AttributeTok{sex =} \FunctionTok{replace\_na}\NormalTok{(sex, }\StringTok{"unknown"}\NormalTok{)}
\NormalTok{  )}
\end{Highlighting}
\end{Shaded}

This approach preserves observations and treats missing values as a
separate category. It may be useful when missingness itself could carry
information, but it also introduces an additional artificial category.

\subsubsection{Option 3: Imputation During Model
Training}\label{option-3-imputation-during-model-training}

\begin{Shaded}
\begin{Highlighting}[]
\NormalTok{model\_data\_option\_3 }\OtherTok{\textless{}{-}}\NormalTok{ model\_data}
\end{Highlighting}
\end{Shaded}

In this approach, missing values are not manually replaced at this
stage. They are handled inside the modelling recipe.

\begin{Shaded}
\begin{Highlighting}[]
\FunctionTok{library}\NormalTok{(tidymodels)}
\NormalTok{model\_recipe }\OtherTok{\textless{}{-}} \FunctionTok{recipe}\NormalTok{(species }\SpecialCharTok{\textasciitilde{}}\NormalTok{ ., }\AttributeTok{data =}\NormalTok{ model\_data\_option\_3) }\SpecialCharTok{|\textgreater{}}
  \FunctionTok{step\_impute\_mode}\NormalTok{(sex) }\SpecialCharTok{|\textgreater{}}
  \FunctionTok{step\_impute\_median}\NormalTok{(}\FunctionTok{all\_numeric\_predictors}\NormalTok{()) }\SpecialCharTok{|\textgreater{}}
  \FunctionTok{step\_dummy}\NormalTok{(}\FunctionTok{all\_nominal\_predictors}\NormalTok{()) }\SpecialCharTok{|\textgreater{}}
  \FunctionTok{step\_normalize}\NormalTok{(}\FunctionTok{all\_numeric\_predictors}\NormalTok{())}
\end{Highlighting}
\end{Shaded}

This strategy keeps preprocessing consistent between training and
testing data and reduces the risk of data leakage.

\subsection{Encoding and Scaling}\label{encoding-and-scaling}

The dataset contains both numerical and categorical variables.
Appropriate preprocessing techniques, including categorical encoding and
numerical scaling where required, were included in the modelling
workflow.

\begin{Shaded}
\begin{Highlighting}[]
\NormalTok{model\_recipe }\OtherTok{\textless{}{-}} \FunctionTok{recipe}\NormalTok{(species }\SpecialCharTok{\textasciitilde{}}\NormalTok{ ., }\AttributeTok{data =}\NormalTok{ model\_data\_option\_3) }\SpecialCharTok{|\textgreater{}}
  \FunctionTok{step\_impute\_mode}\NormalTok{(}\FunctionTok{all\_nominal\_predictors}\NormalTok{()) }\SpecialCharTok{|\textgreater{}}
  \FunctionTok{step\_impute\_median}\NormalTok{(}\FunctionTok{all\_numeric\_predictors}\NormalTok{()) }\SpecialCharTok{|\textgreater{}}
  \FunctionTok{step\_dummy}\NormalTok{(}\FunctionTok{all\_nominal\_predictors}\NormalTok{()) }\SpecialCharTok{|\textgreater{}}
  \FunctionTok{step\_normalize}\NormalTok{(}\FunctionTok{all\_numeric\_predictors}\NormalTok{())}
\end{Highlighting}
\end{Shaded}

Categorical variables were encoded into numerical form using dummy
variables. Numerical variables were normalized to make them comparable
in scale, which is especially important for models sensitive to variable
magnitude.

\subsection{Training and Testing Sets}\label{training-and-testing-sets}

To evaluate model performance objectively, the dataset was divided into
separate training and testing subsets. The training set was used for
model development, while the testing set was reserved for final
performance evaluation.

\begin{Shaded}
\begin{Highlighting}[]
\FunctionTok{set.seed}\NormalTok{(}\DecValTok{123}\NormalTok{)}

\NormalTok{data\_split }\OtherTok{\textless{}{-}} \FunctionTok{initial\_split}\NormalTok{(}
\NormalTok{  model\_data\_option\_3,}
  \AttributeTok{prop =} \FloatTok{0.8}\NormalTok{,}
  \AttributeTok{strata =}\NormalTok{ species}
\NormalTok{)}

\NormalTok{train\_data }\OtherTok{\textless{}{-}} \FunctionTok{training}\NormalTok{(data\_split)}
\NormalTok{test\_data  }\OtherTok{\textless{}{-}} \FunctionTok{testing}\NormalTok{(data\_split)}
\end{Highlighting}
\end{Shaded}

\begin{Shaded}
\begin{Highlighting}[]
\NormalTok{model\_recipe }\OtherTok{\textless{}{-}} \FunctionTok{recipe}\NormalTok{(species }\SpecialCharTok{\textasciitilde{}}\NormalTok{ ., }\AttributeTok{data =}\NormalTok{ train\_data) }\SpecialCharTok{|\textgreater{}}
  \FunctionTok{step\_impute\_mode}\NormalTok{(}\FunctionTok{all\_nominal\_predictors}\NormalTok{()) }\SpecialCharTok{|\textgreater{}}
  \FunctionTok{step\_impute\_median}\NormalTok{(}\FunctionTok{all\_numeric\_predictors}\NormalTok{()) }\SpecialCharTok{|\textgreater{}}
  \FunctionTok{step\_dummy}\NormalTok{(}\FunctionTok{all\_nominal\_predictors}\NormalTok{()) }\SpecialCharTok{|\textgreater{}}
  \FunctionTok{step\_normalize}\NormalTok{(}\FunctionTok{all\_numeric\_predictors}\NormalTok{())}
\end{Highlighting}
\end{Shaded}

The split was stratified by the target variable to preserve the
distribution of penguin species in both subsets.

\subsection{Prepared Data for Model 1}\label{prepared-data-for-model-1}

The prepared recipe can now be used as part of the modelling workflow.
This ensures that all preprocessing steps are applied consistently
during model training and evaluation.

\begin{Shaded}
\begin{Highlighting}[]
\NormalTok{model\_1\_spec }\OtherTok{\textless{}{-}} \FunctionTok{decision\_tree}\NormalTok{() }\SpecialCharTok{|\textgreater{}}
  \FunctionTok{set\_engine}\NormalTok{(}\StringTok{"rpart"}\NormalTok{) }\SpecialCharTok{|\textgreater{}}
  \FunctionTok{set\_mode}\NormalTok{(}\StringTok{"classification"}\NormalTok{)}

\NormalTok{model\_1\_workflow }\OtherTok{\textless{}{-}} \FunctionTok{workflow}\NormalTok{() }\SpecialCharTok{|\textgreater{}}
  \FunctionTok{add\_recipe}\NormalTok{(model\_recipe) }\SpecialCharTok{|\textgreater{}}
  \FunctionTok{add\_model}\NormalTok{(model\_1\_spec)}
\end{Highlighting}
\end{Shaded}

\subsection{Summary}\label{summary}

The dataset was cleaned and prepared for predictive modelling. The
resulting data contains consistent variable definitions and provides a
suitable basis for building and evaluating classification models.

The final preparation workflow includes:

\begin{itemize}
\tightlist
\item
  removal of non-informative variables,
\item
  standardization of categorical labels,
\item
  treatment of implausible values,
\item
  imputation of remaining missing values,
\item
  encoding of categorical predictors,
\item
  normalization of numerical predictors,
\item
  division into training and testing subsets.
\end{itemize}

These steps provide a structured and reproducible foundation for the
modelling stage.

\subsection{Check \#2: Data Validation}\label{check-2-data-validation}

After the initial cleaning step, the dataset was checked again to verify
whether implausible values had been removed successfully. This
validation revealed that some unrealistic observations were still
present and required additional filtering.

\begin{Shaded}
\begin{Highlighting}[]
\FunctionTok{summary}\NormalTok{(model\_data)}
\end{Highlighting}
\end{Shaded}

\begin{verbatim}
##    species          bill_length_mm  bill_depth_mm   flipper_length_mm
##  Length:429         Min.   :31.20   Min.   :12.60   Min.   :  1.0    
##  Class :character   1st Qu.:39.50   1st Qu.:15.70   1st Qu.:189.0    
##  Mode  :character   Median :43.60   Median :17.70   Median :196.0    
##                     Mean   :43.71   Mean   :17.27   Mean   :199.5    
##                     3rd Qu.:48.10   3rd Qu.:18.70   3rd Qu.:213.0    
##                     Max.   :58.30   Max.   :21.70   Max.   :232.0    
##                                     NA's   :19                       
##   body_mass_g      island              sex                 year     
##  Min.   :-500   Length:429         Length:429         Min.   :2007  
##  1st Qu.:3562   Class :character   Class :character   1st Qu.:2007  
##  Median :4039   Mode  :character   Mode  :character   Median :2008  
##  Mean   :4118                                         Mean   :2008  
##  3rd Qu.:4795                                         3rd Qu.:2009  
##  Max.   :6843                                         Max.   :2099  
## 
\end{verbatim}

There are still strange data: - Min body\_mass\_g = -500 - Min
flipper\_length\_mm = 1 - Max year = 2099 So:

\begin{Shaded}
\begin{Highlighting}[]
\NormalTok{model\_data }\OtherTok{\textless{}{-}}\NormalTok{ model\_data }\SpecialCharTok{|\textgreater{}}
  \FunctionTok{filter}\NormalTok{(}
    \FunctionTok{between}\NormalTok{(bill\_length\_mm, }\DecValTok{20}\NormalTok{, }\DecValTok{70}\NormalTok{) }\SpecialCharTok{|} \FunctionTok{is.na}\NormalTok{(bill\_length\_mm),}
    \FunctionTok{between}\NormalTok{(flipper\_length\_mm, }\DecValTok{150}\NormalTok{, }\DecValTok{250}\NormalTok{) }\SpecialCharTok{|} \FunctionTok{is.na}\NormalTok{(flipper\_length\_mm),}
    \FunctionTok{between}\NormalTok{(body\_mass\_g, }\DecValTok{2000}\NormalTok{, }\DecValTok{8000}\NormalTok{) }\SpecialCharTok{|} \FunctionTok{is.na}\NormalTok{(body\_mass\_g),}
    \FunctionTok{between}\NormalTok{(year, }\DecValTok{2007}\NormalTok{, }\DecValTok{2009}\NormalTok{) }\SpecialCharTok{|} \FunctionTok{is.na}\NormalTok{(year)}
\NormalTok{  )}
\end{Highlighting}
\end{Shaded}

\begin{Shaded}
\begin{Highlighting}[]
\FunctionTok{summary}\NormalTok{(model\_data)}
\end{Highlighting}
\end{Shaded}

\begin{verbatim}
##    species          bill_length_mm  bill_depth_mm   flipper_length_mm
##  Length:419         Min.   :31.20   Min.   :12.60   Min.   :172.0    
##  Class :character   1st Qu.:39.45   1st Qu.:15.60   1st Qu.:189.0    
##  Mode  :character   Median :43.50   Median :17.70   Median :196.0    
##                     Mean   :43.67   Mean   :17.27   Mean   :199.8    
##                     3rd Qu.:48.10   3rd Qu.:18.80   3rd Qu.:213.0    
##                     Max.   :58.30   Max.   :21.70   Max.   :232.0    
##                                     NA's   :19                       
##   body_mass_g      island              sex                 year     
##  Min.   :2399   Length:419         Length:419         Min.   :2007  
##  1st Qu.:3596   Class :character   Class :character   1st Qu.:2007  
##  Median :4061   Mode  :character   Mode  :character   Median :2008  
##  Mean   :4187                                         Mean   :2008  
##  3rd Qu.:4811                                         3rd Qu.:2009  
##  Max.   :6843                                         Max.   :2009  
## 
\end{verbatim}

he additional validation step confirmed that extreme values remaining
after the first cleaning pass were removed. This illustrates that data
preparation is an iterative process: cleaning decisions should be
followed by checks to verify whether the dataset is actually ready for
modelling.

\subsection{Models}\label{models}

Before fitting classification models, the target variable and
categorical predictors were converted to factors. This ensures that the
modelling workflow treats the task as classification rather than
regression.

\begin{Shaded}
\begin{Highlighting}[]
\NormalTok{model\_data\_option\_3 }\OtherTok{\textless{}{-}}\NormalTok{ model\_data\_option\_3 }\SpecialCharTok{|\textgreater{}}
  \FunctionTok{mutate}\NormalTok{(}
    \AttributeTok{species =} \FunctionTok{as.factor}\NormalTok{(species),}
    \AttributeTok{island =} \FunctionTok{as.factor}\NormalTok{(island),}
    \AttributeTok{sex =} \FunctionTok{as.factor}\NormalTok{(sex)}
\NormalTok{  )}
\end{Highlighting}
\end{Shaded}

\begin{Shaded}
\begin{Highlighting}[]
\FunctionTok{set.seed}\NormalTok{(}\DecValTok{123}\NormalTok{)}

\NormalTok{data\_split }\OtherTok{\textless{}{-}} \FunctionTok{initial\_split}\NormalTok{(}
\NormalTok{  model\_data\_option\_3,}
  \AttributeTok{prop =} \FloatTok{0.8}\NormalTok{,}
  \AttributeTok{strata =}\NormalTok{ species}
\NormalTok{)}

\NormalTok{train\_data }\OtherTok{\textless{}{-}} \FunctionTok{training}\NormalTok{(data\_split)}
\NormalTok{test\_data  }\OtherTok{\textless{}{-}} \FunctionTok{testing}\NormalTok{(data\_split)}
\end{Highlighting}
\end{Shaded}

\begin{Shaded}
\begin{Highlighting}[]
\NormalTok{model\_recipe }\OtherTok{\textless{}{-}} \FunctionTok{recipe}\NormalTok{(species }\SpecialCharTok{\textasciitilde{}}\NormalTok{ ., }\AttributeTok{data =}\NormalTok{ train\_data) }\SpecialCharTok{|\textgreater{}}
  \FunctionTok{step\_impute\_mode}\NormalTok{(}\FunctionTok{all\_nominal\_predictors}\NormalTok{()) }\SpecialCharTok{|\textgreater{}}
  \FunctionTok{step\_impute\_median}\NormalTok{(}\FunctionTok{all\_numeric\_predictors}\NormalTok{()) }\SpecialCharTok{|\textgreater{}}
  \FunctionTok{step\_dummy}\NormalTok{(}\FunctionTok{all\_nominal\_predictors}\NormalTok{())}
\end{Highlighting}
\end{Shaded}

\begin{Shaded}
\begin{Highlighting}[]
\NormalTok{classification\_metrics }\OtherTok{\textless{}{-}} \FunctionTok{metric\_set}\NormalTok{(}
\NormalTok{  accuracy,}
\NormalTok{  kap}
\NormalTok{)}
\end{Highlighting}
\end{Shaded}

\subsubsection{Model 1}\label{model-1}

\paragraph{Model Description}\label{model-description}

A Decision Tree classifier was selected as the first model. This model
is easy to interpret and provides a useful baseline for classification
tasks. Decision trees split the data into smaller groups based on
predictor variables and generate classification rules that can be
inspected and explained.

This model was selected because it is transparent, relatively simple,
and suitable as an initial reference model for predicting penguin
species.

\paragraph{Implementation}\label{implementation}

\begin{Shaded}
\begin{Highlighting}[]
\NormalTok{dt\_model }\OtherTok{\textless{}{-}} \FunctionTok{decision\_tree}\NormalTok{() }\SpecialCharTok{|\textgreater{}}
  \FunctionTok{set\_engine}\NormalTok{(}\StringTok{"rpart"}\NormalTok{) }\SpecialCharTok{|\textgreater{}}
  \FunctionTok{set\_mode}\NormalTok{(}\StringTok{"classification"}\NormalTok{)}

\NormalTok{dt\_workflow }\OtherTok{\textless{}{-}} \FunctionTok{workflow}\NormalTok{() }\SpecialCharTok{|\textgreater{}}
  \FunctionTok{add\_recipe}\NormalTok{(model\_recipe) }\SpecialCharTok{|\textgreater{}}
  \FunctionTok{add\_model}\NormalTok{(dt\_model)}

\NormalTok{dt\_fit }\OtherTok{\textless{}{-}} \FunctionTok{fit}\NormalTok{(}
\NormalTok{  dt\_workflow,}
  \AttributeTok{data =}\NormalTok{ train\_data}
\NormalTok{)}
\end{Highlighting}
\end{Shaded}

Which gives:

\begin{Shaded}
\begin{Highlighting}[]
\FunctionTok{library}\NormalTok{(rpart.plot)}

\NormalTok{dt\_tree }\OtherTok{\textless{}{-}} \FunctionTok{extract\_fit\_engine}\NormalTok{(dt\_fit)}

\FunctionTok{rpart.plot}\NormalTok{(}
\NormalTok{  dt\_tree,}
  \AttributeTok{type =} \DecValTok{2}\NormalTok{,}
  \AttributeTok{extra =} \DecValTok{104}\NormalTok{,}
  \AttributeTok{fallen.leaves =} \ConstantTok{TRUE}\NormalTok{,}
  \AttributeTok{main =} \StringTok{"Decision Tree for Penguin Species Classification"}
\NormalTok{)}
\end{Highlighting}
\end{Shaded}

\pandocbounded{\includegraphics[keepaspectratio]{example-report_files/figures/unnamed-chunk-45-1.pdf}}
The decision tree visualization shows the main decision rules used by
the model to classify penguin species. This makes the model easier to
interpret than more complex ensemble methods. However, a single tree may
be less stable and less accurate than models such as Random Forest.
\#\#\#\# Results

\begin{Shaded}
\begin{Highlighting}[]
\NormalTok{dt\_predictions }\OtherTok{\textless{}{-}} \FunctionTok{predict}\NormalTok{(}
\NormalTok{  dt\_fit,}
\NormalTok{  test\_data}
\NormalTok{) }\SpecialCharTok{|\textgreater{}}
  \FunctionTok{bind\_cols}\NormalTok{(test\_data)}

\NormalTok{dt\_metrics }\OtherTok{\textless{}{-}}\NormalTok{ dt\_predictions }\SpecialCharTok{|\textgreater{}}
  \FunctionTok{classification\_metrics}\NormalTok{(}
    \AttributeTok{truth =}\NormalTok{ species,}
    \AttributeTok{estimate =}\NormalTok{ .pred\_class}
\NormalTok{  )}

\NormalTok{dt\_metrics}
\end{Highlighting}
\end{Shaded}

\begin{verbatim}
## # A tibble: 2 x 3
##   .metric  .estimator .estimate
##   <chr>    <chr>          <dbl>
## 1 accuracy multiclass     0.966
## 2 kap      multiclass     0.946
\end{verbatim}

\begin{Shaded}
\begin{Highlighting}[]
\NormalTok{dt\_conf\_mat }\OtherTok{\textless{}{-}}\NormalTok{ dt\_predictions }\SpecialCharTok{|\textgreater{}}
  \FunctionTok{conf\_mat}\NormalTok{(}
    \AttributeTok{truth =}\NormalTok{ species,}
    \AttributeTok{estimate =}\NormalTok{ .pred\_class}
\NormalTok{  )}

\NormalTok{dt\_conf\_mat}
\end{Highlighting}
\end{Shaded}

\begin{verbatim}
##            Truth
## Prediction  Adelie Chinstrap Gentoo
##   Adelie        37         0      0
##   Chinstrap      3        17      0
##   Gentoo         0         0     30
\end{verbatim}

\paragraph{Interpretation}\label{interpretation-4}

The Decision Tree model provides a baseline classification result. Its
performance metrics show how well a relatively simple and interpretable
model can distinguish between penguin species using the selected
predictor variables.

The confusion matrix can be used to identify which species are
classified correctly and where misclassifications occur.

\subsubsection{Model 2}\label{model-2}

\paragraph{Model Description}\label{model-description-1}

A Random Forest classifier was selected as the second model. Random
Forest extends the idea of a single decision tree by combining multiple
trees into an ensemble model. This approach often improves predictive
performance and reduces the risk of overfitting.

This model was selected because it is robust, works well with mixed
predictor variables, and often provides strong performance in
classification problems.

\paragraph{Implementation}\label{implementation-1}

\begin{Shaded}
\begin{Highlighting}[]
\FunctionTok{library}\NormalTok{(ranger)}
\NormalTok{rf\_model }\OtherTok{\textless{}{-}} \FunctionTok{rand\_forest}\NormalTok{(}
  \AttributeTok{trees =} \DecValTok{500}
\NormalTok{) }\SpecialCharTok{|\textgreater{}}
  \FunctionTok{set\_engine}\NormalTok{(}\StringTok{"ranger"}\NormalTok{, }\AttributeTok{importance =} \StringTok{"impurity"}\NormalTok{) }\SpecialCharTok{|\textgreater{}}
  \FunctionTok{set\_mode}\NormalTok{(}\StringTok{"classification"}\NormalTok{)}

\NormalTok{rf\_workflow }\OtherTok{\textless{}{-}} \FunctionTok{workflow}\NormalTok{() }\SpecialCharTok{|\textgreater{}}
  \FunctionTok{add\_recipe}\NormalTok{(model\_recipe) }\SpecialCharTok{|\textgreater{}}
  \FunctionTok{add\_model}\NormalTok{(rf\_model)}

\NormalTok{rf\_fit }\OtherTok{\textless{}{-}} \FunctionTok{fit}\NormalTok{(}
\NormalTok{  rf\_workflow,}
  \AttributeTok{data =}\NormalTok{ train\_data}
\NormalTok{)}
\end{Highlighting}
\end{Shaded}

Checking variable importance:

\begin{Shaded}
\begin{Highlighting}[]
\FunctionTok{library}\NormalTok{(vip)}

\FunctionTok{vip}\NormalTok{(}
  \FunctionTok{extract\_fit\_parsnip}\NormalTok{(rf\_fit)}\SpecialCharTok{$}\NormalTok{fit,}
  \AttributeTok{num\_features =} \DecValTok{10}
\NormalTok{)}
\end{Highlighting}
\end{Shaded}

\pandocbounded{\includegraphics[keepaspectratio]{example-report_files/figures/unnamed-chunk-49-1.pdf}}
Variable importance indicates which predictors contributed most to the
classification process. The plot helps identify variables that are most
influential when distinguishing between penguin species.

With \texttt{ggplot} we can confirm:

\begin{Shaded}
\begin{Highlighting}[]
\NormalTok{rf\_engine }\OtherTok{\textless{}{-}} \FunctionTok{extract\_fit\_engine}\NormalTok{(rf\_fit)}

\NormalTok{importance\_df }\OtherTok{\textless{}{-}} \FunctionTok{tibble}\NormalTok{(}
  \AttributeTok{Variable =} \FunctionTok{names}\NormalTok{(rf\_engine}\SpecialCharTok{$}\NormalTok{variable.importance),}
  \AttributeTok{Importance =}\NormalTok{ rf\_engine}\SpecialCharTok{$}\NormalTok{variable.importance}
\NormalTok{)}

\NormalTok{importance\_df }\SpecialCharTok{|\textgreater{}}
  \FunctionTok{arrange}\NormalTok{(}\FunctionTok{desc}\NormalTok{(Importance)) }\SpecialCharTok{|\textgreater{}}
  \FunctionTok{ggplot}\NormalTok{(}\FunctionTok{aes}\NormalTok{(}
    \AttributeTok{x =} \FunctionTok{reorder}\NormalTok{(Variable, Importance),}
    \AttributeTok{y =}\NormalTok{ Importance}
\NormalTok{  )) }\SpecialCharTok{+}
  \FunctionTok{geom\_col}\NormalTok{() }\SpecialCharTok{+}
  \FunctionTok{coord\_flip}\NormalTok{() }\SpecialCharTok{+}
  \FunctionTok{labs}\NormalTok{(}
    \AttributeTok{title =} \StringTok{"Random Forest Variable Importance"}\NormalTok{,}
    \AttributeTok{x =} \StringTok{""}\NormalTok{,}
    \AttributeTok{y =} \StringTok{"Importance"}
\NormalTok{  )}
\end{Highlighting}
\end{Shaded}

\pandocbounded{\includegraphics[keepaspectratio]{example-report_files/figures/unnamed-chunk-50-1.pdf}}

\paragraph{Results}\label{results}

\begin{Shaded}
\begin{Highlighting}[]
\NormalTok{rf\_predictions }\OtherTok{\textless{}{-}} \FunctionTok{predict}\NormalTok{(}
\NormalTok{  rf\_fit,}
\NormalTok{  test\_data}
\NormalTok{) }\SpecialCharTok{|\textgreater{}}
  \FunctionTok{bind\_cols}\NormalTok{(test\_data)}

\NormalTok{rf\_metrics }\OtherTok{\textless{}{-}}\NormalTok{ rf\_predictions }\SpecialCharTok{|\textgreater{}}
  \FunctionTok{classification\_metrics}\NormalTok{(}
    \AttributeTok{truth =}\NormalTok{ species,}
    \AttributeTok{estimate =}\NormalTok{ .pred\_class}
\NormalTok{  )}

\NormalTok{rf\_metrics}
\end{Highlighting}
\end{Shaded}

\begin{verbatim}
## # A tibble: 2 x 3
##   .metric  .estimator .estimate
##   <chr>    <chr>          <dbl>
## 1 accuracy multiclass     0.977
## 2 kap      multiclass     0.964
\end{verbatim}

\begin{Shaded}
\begin{Highlighting}[]
\NormalTok{rf\_conf\_mat }\OtherTok{\textless{}{-}}\NormalTok{ rf\_predictions }\SpecialCharTok{|\textgreater{}}
  \FunctionTok{conf\_mat}\NormalTok{(}
    \AttributeTok{truth =}\NormalTok{ species,}
    \AttributeTok{estimate =}\NormalTok{ .pred\_class}
\NormalTok{  )}

\NormalTok{rf\_conf\_mat}
\end{Highlighting}
\end{Shaded}

\begin{verbatim}
##            Truth
## Prediction  Adelie Chinstrap Gentoo
##   Adelie        38         0      0
##   Chinstrap      2        17      0
##   Gentoo         0         0     30
\end{verbatim}

\paragraph{Interpretation}\label{interpretation-5}

The Random Forest model was evaluated using the testing dataset. The
obtained metrics indicate how effectively the model predicts penguin
species based on the selected variables.

The results can be compared with the Decision Tree model to assess
whether the ensemble approach improves predictive performance.

\subsection{Model Comparison}\label{model-comparison}

\subsubsection{Results}\label{results-1}

\begin{Shaded}
\begin{Highlighting}[]
\NormalTok{model\_comparison }\OtherTok{\textless{}{-}} \FunctionTok{bind\_rows}\NormalTok{(}
\NormalTok{  dt\_metrics }\SpecialCharTok{|\textgreater{}} \FunctionTok{mutate}\NormalTok{(}\AttributeTok{model =} \StringTok{"Decision Tree"}\NormalTok{),}
\NormalTok{  rf\_metrics }\SpecialCharTok{|\textgreater{}} \FunctionTok{mutate}\NormalTok{(}\AttributeTok{model =} \StringTok{"Random Forest"}\NormalTok{)}
\NormalTok{) }\SpecialCharTok{|\textgreater{}}
  \FunctionTok{select}\NormalTok{(model, .metric, .estimate)}

\NormalTok{model\_comparison}
\end{Highlighting}
\end{Shaded}

\begin{verbatim}
## # A tibble: 4 x 3
##   model         .metric  .estimate
##   <chr>         <chr>        <dbl>
## 1 Decision Tree accuracy     0.966
## 2 Decision Tree kap          0.946
## 3 Random Forest accuracy     0.977
## 4 Random Forest kap          0.964
\end{verbatim}

\begin{Shaded}
\begin{Highlighting}[]
\NormalTok{model\_comparison }\SpecialCharTok{|\textgreater{}}
  \FunctionTok{ggplot}\NormalTok{(}\FunctionTok{aes}\NormalTok{(}\AttributeTok{x =}\NormalTok{ model, }\AttributeTok{y =}\NormalTok{ .estimate, }\AttributeTok{fill =}\NormalTok{ .metric)) }\SpecialCharTok{+}
  \FunctionTok{geom\_col}\NormalTok{(}\AttributeTok{position =} \StringTok{"dodge"}\NormalTok{) }\SpecialCharTok{+}
  \FunctionTok{labs}\NormalTok{(}
    \AttributeTok{title =} \StringTok{"Model comparison"}\NormalTok{,}
    \AttributeTok{x =} \StringTok{"Model"}\NormalTok{,}
    \AttributeTok{y =} \StringTok{"Metric value"}
\NormalTok{  )}
\end{Highlighting}
\end{Shaded}

\pandocbounded{\includegraphics[keepaspectratio]{example-report_files/figures/unnamed-chunk-54-1.pdf}}

\subsubsection{Interpretation}\label{interpretation-6}

Both models were evaluated using the same testing dataset. The
comparison allows assessment of whether the more complex Random Forest
model provides better predictive performance than the simpler Decision
Tree model.

The final model selection should consider both predictive performance
and interpretability.

\subsection{Discussion}\label{discussion}

Two classification models were developed and evaluated: a Decision Tree
and a Random Forest.

The Decision Tree model offers a high level of interpretability. The
resulting tree structure can be visualized and easily understood, making
it useful for explaining classification decisions. However, a single
tree may be sensitive to variations in the training data and can be
prone to overfitting.

The Random Forest model addresses some of these limitations by combining
multiple decision trees into an ensemble. This approach generally
improves predictive performance and robustness while reducing the impact
of individual tree errors. The trade-off is reduced interpretability, as
the model's decision-making process becomes more complex.

Overall, both models demonstrate how machine learning techniques can be
applied to classification problems. The comparison highlights the common
trade-off between model interpretability and predictive performance.

\subsection{Limitations}\label{limitations}

Several limitations should be considered when interpreting the results.

First, the dataset required extensive cleaning due to missing values,
inconsistent categorical labels, duplicate records, and implausible
observations. The quality of the final models therefore depends on the
assumptions made during the data preparation stage.

Second, the analysis was conducted using a single train-test split.
Different random splits or cross-validation procedures could produce
slightly different performance estimates.

Third, the project focused on a limited number of classification
algorithms. Other approaches, such as k-Nearest Neighbours, Support
Vector Machines, Gradient Boosting, or Neural Networks, may yield
different results.

Finally, the dataset represents a specific collection of penguin
observations. Consequently, conclusions should be interpreted within the
context of this dataset and should not be generalized beyond the
available data without additional validation.

\section{7. Recommendations}\label{recommendations}

Provide recommendations based on the results of the analysis.

\section{8. Conclusions}\label{conclusions}

The objective of this project was to develop and evaluate machine
learning models capable of predicting penguin species based on
morphological measurements and selected categorical attributes.

The project objectives were successfully achieved. Data quality issues,
including missing values, inconsistent category labels, duplicate
observations, and implausible values, were identified and addressed
during the data preparation stage. This improved the overall consistency
and usability of the dataset.

A meaningful exploratory data analysis was performed, providing insights
into variable distributions, relationships between predictors,
correlations, and potential outliers. These analyses supported
subsequent data preparation and modelling decisions.

Two classification models, a Decision Tree and a Random Forest, were
developed and evaluated using the prepared dataset. Model performance
was assessed using objective evaluation metrics and confusion matrices,
allowing a direct comparison of the two approaches.

The results demonstrate that machine learning techniques can effectively
classify penguin species using the available predictor variables. While
the Decision Tree offered greater interpretability, the Random Forest
generally provided a more robust classification approach.

Overall, the success criteria defined at the beginning of the project
were met:

\begin{itemize}
\tightlist
\item
  data quality issues were identified and addressed;
\item
  meaningful exploratory analysis was performed;
\item
  two classification models were developed;
\item
  model performance was evaluated and compared;
\item
  conclusions were supported by analytical evidence obtained throughout
  the CRISP-DM process.
\end{itemize}

The project illustrates the importance of systematic data preparation,
exploratory analysis, and model evaluation when developing predictive
solutions based on real-world data.

\subsection{Summary}\label{summary-1}

This project followed the CRISP-DM methodology to develop and evaluate
machine learning models for penguin species classification.

The analysis included data understanding, data cleaning, exploratory
data analysis, feature preparation, model development, and model
evaluation. Data quality issues were identified and addressed prior to
modelling. Two classification algorithms, a Decision Tree and a Random
Forest, were trained and evaluated using a dedicated testing dataset.

The results demonstrate the importance of systematic data preparation
and model evaluation when building predictive solutions based on
real-world data.

\subsection{Lessons Learned}\label{lessons-learned}

Several important lessons emerged from this analysis.

First, data quality has a substantial impact on the modelling process.
Missing values, inconsistent labels, duplicate observations, and
implausible values required attention before reliable models could be
developed.

Second, exploratory data analysis provided valuable insights into the
structure of the dataset and helped guide data preparation decisions.
Visualisation proved particularly useful for identifying unusual
observations and understanding relationships between variables.

Third, different machine learning models offer different advantages.
Simpler models such as Decision Trees provide interpretability, while
ensemble methods such as Random Forests often deliver improved
predictive performance.

Finally, model evaluation should extend beyond a single metric and
include multiple perspectives, such as confusion matrices and
comparative performance analysis.

\subsection{Future Work}\label{future-work}

Several extensions could be explored in future iterations of the
project.

Additional classification algorithms, such as Support Vector Machines,
k-Nearest Neighbours, Gradient Boosting, or Neural Networks, could be
evaluated and compared with the models developed in this study.

More advanced feature engineering techniques could also be investigated
to improve predictive performance. Cross-validation procedures could
provide more robust estimates of model generalisation performance.

Finally, the analysis could be extended by incorporating additional
biological or environmental variables, allowing more comprehensive
modelling of penguin characteristics and species differentiation.

\section{References}\label{references}

List all external sources used.

\section{Appendix A (Optional)}\label{appendix-a-optional}

Additional tables, figures, code snippets, or supplementary material.

\section{Appendix B -Explanatory Data Analysis
Report}\label{appendix-b--explanatory-data-analysis-report}

💡 Accessibility and Visual Design

The visualisations in this report were designed with readability and
accessibility in mind. Colour palettes and formatting choices were
selected to ensure that charts remain interpretable for a wide audience,
including individuals with common forms of colour vision deficiency.

In addition to colour, visual distinctions such as point shapes, labels,
and layout were used to improve clarity. These design choices help
ensure that the presented insights remain understandable when viewed on
different screens, projectors, or printed materials.

Effective data communication depends not only on accurate analysis but
also on the ability to present results in a clear, accessible, and
audience-friendly manner.

\subsection{Understanding Penguin
Characteristics}\label{understanding-penguin-characteristics}

The purpose of this report is not to build predictive models, but to
communicate selected insights discovered in the penguin dataset. While
exploratory analysis focuses on discovering patterns, explanatory
analysis focuses on presenting findings in a clear and meaningful way
for an audience.

The dataset contains measurements of penguins belonging to three
species: Adelie, Chinstrap, and Gentoo. Several physical characteristics
were recorded, including bill dimensions, flipper length, body mass,
sex, and island of origin.

\subsection{Insight 1: Penguin Species Differ in Physical
Characteristics}\label{insight-1-penguin-species-differ-in-physical-characteristics}

One of the most visible patterns in the dataset is the difference in
body measurements between species.

When comparing bill dimensions and flipper length, the three species
form distinct groups. Gentoo penguins generally have longer flippers and
greater body mass, while Adelie penguins tend to be smaller. Chinstrap
penguins often occupy an intermediate position.

These differences suggest that morphological measurements contain useful
information for distinguishing between species.

\begin{Shaded}
\begin{Highlighting}[]
\FunctionTok{ggplot}\NormalTok{(}
\NormalTok{  model\_data,}
  \FunctionTok{aes}\NormalTok{(}
    \AttributeTok{x =}\NormalTok{ bill\_length\_mm,}
    \AttributeTok{y =}\NormalTok{ flipper\_length\_mm,}
    \AttributeTok{color =}\NormalTok{ species,}
    \AttributeTok{shape =}\NormalTok{ species}
\NormalTok{  )}
\NormalTok{) }\SpecialCharTok{+}
  \FunctionTok{geom\_point}\NormalTok{(}
    \AttributeTok{size =} \DecValTok{2}\NormalTok{,}
    \AttributeTok{alpha =} \FloatTok{0.8}
\NormalTok{  ) }\SpecialCharTok{+}
  \FunctionTok{scale\_color\_manual}\NormalTok{(}
    \AttributeTok{values =} \FunctionTok{c}\NormalTok{(}
      \StringTok{"Adelie"} \OtherTok{=} \StringTok{"\#0072B2"}\NormalTok{,}
      \StringTok{"Chinstrap"} \OtherTok{=} \StringTok{"\#E69F00"}\NormalTok{,}
      \StringTok{"Gentoo"} \OtherTok{=} \StringTok{"\#009E73"}
\NormalTok{    )}
\NormalTok{  ) }\SpecialCharTok{+}
  \FunctionTok{theme\_minimal}\NormalTok{()}
\end{Highlighting}
\end{Shaded}

\pandocbounded{\includegraphics[keepaspectratio]{example-report_files/figures/unnamed-chunk-55-1.pdf}}
The three penguin species occupy different regions of the feature space.
This suggests that morphological measurements contain sufficient
information to distinguish species and explains why classification
models achieve good predictive performance.

\subsection{Insight 2: Species Distribution Across
Islands}\label{insight-2-species-distribution-across-islands}

Species distribution was examined across islands to explore whether
penguin species occur uniformly throughout the study area. Understanding
geographic distribution patterns can provide additional biological
context and help explain why location-related variables may be useful
for classification

\begin{Shaded}
\begin{Highlighting}[]
\FunctionTok{ggplot}\NormalTok{(}
\NormalTok{  data,}
  \FunctionTok{aes}\NormalTok{(}
    \AttributeTok{x =}\NormalTok{ island,}
    \AttributeTok{fill =}\NormalTok{ species}
\NormalTok{  )}
\NormalTok{) }\SpecialCharTok{+}
  \FunctionTok{geom\_bar}\NormalTok{(}
    \AttributeTok{position =} \StringTok{"fill"}\NormalTok{,}
    \AttributeTok{color =} \StringTok{"white"}\NormalTok{,}
    \AttributeTok{linewidth =} \FloatTok{0.3}
\NormalTok{  ) }\SpecialCharTok{+}
  \FunctionTok{scale\_fill\_manual}\NormalTok{(}
    \AttributeTok{values =} \FunctionTok{c}\NormalTok{(}
      \StringTok{"Adelie"} \OtherTok{=} \StringTok{"\#0072B2"}\NormalTok{,}
      \StringTok{"Chinstrap"} \OtherTok{=} \StringTok{"\#E69F00"}\NormalTok{,}
      \StringTok{"Gentoo"} \OtherTok{=} \StringTok{"\#009E73"}
\NormalTok{    )}
\NormalTok{  ) }\SpecialCharTok{+}
  \FunctionTok{scale\_y\_continuous}\NormalTok{(}
    \AttributeTok{labels =}\NormalTok{ scales}\SpecialCharTok{::}\NormalTok{percent}
\NormalTok{  ) }\SpecialCharTok{+}
  \FunctionTok{theme\_minimal}\NormalTok{() }\SpecialCharTok{+}
  \FunctionTok{labs}\NormalTok{(}
    \AttributeTok{title =} \StringTok{"Species Composition Differs Across Islands"}\NormalTok{,}
    \AttributeTok{subtitle =} \StringTok{"Each island exhibits a distinct species profile"}\NormalTok{,}
    \AttributeTok{x =} \StringTok{"Island"}\NormalTok{,}
    \AttributeTok{y =} \StringTok{"Percentage of Observations"}\NormalTok{,}
    \AttributeTok{fill =} \StringTok{"Species"}
\NormalTok{  )}
\end{Highlighting}
\end{Shaded}

\pandocbounded{\includegraphics[keepaspectratio]{example-report_files/figures/unnamed-chunk-56-1.pdf}}

One of the clearest patterns in the dataset is that penguin species are
not evenly distributed across islands. Certain islands are associated
with particular species, creating distinct geographic signatures. This
observation demonstrates that location alone can provide valuable clues
about species identity and helps explain why the island variable
contributes to classification performance

\subsection{Insight 3: Data Quality Directly Affects Analysis
Quality}\label{insight-3-data-quality-directly-affects-analysis-quality}

The dataset contained several data quality issues, including:

\begin{itemize}
\tightlist
\item
  missing values,
\item
  duplicate observations,
\item
  inconsistent category labels,
\item
  implausible numerical values.
\end{itemize}

These problems required cleaning before reliable analysis could be
performed.

\subsubsection{Why It Matters}\label{why-it-matters}

Poor data quality can distort visualisations, bias statistical
summaries, and reduce model performance. Data preparation is therefore a
critical step in any data mining project.

\subsection{Insight 4: Visualisation Reveals Patterns Hidden in Summary
Statistics}\label{insight-4-visualisation-reveals-patterns-hidden-in-summary-statistics}

Several relationships became visible only after visual exploration of
the data.

Histograms revealed unusual observations and potential outliers.
Scatterplots highlighted relationships between variables. Correlation
analysis quantified some of these relationships, while boxplots
illustrated differences between groups.

\subsubsection{Why It Matters}\label{why-it-matters-1}

Numerical summaries alone rarely provide a complete understanding of a
dataset. Visualisation remains one of the most effective tools for
communicating analytical findings.

\subsection{Key Takeaways}\label{key-takeaways}

The analysis demonstrates that:

\begin{itemize}
\tightlist
\item
  penguin species exhibit distinct physical characteristics;
\item
  body mass differs between male and female penguins;
\item
  data quality strongly influences analytical outcomes;
\item
  visualisation is essential for discovering and communicating insights.
\end{itemize}

\end{document}
